\section{Aperture-dependent minimax tomography with coherent-state tests}
\label{sec:coherent-aperture}
This section gives matching deterministic and statistical rates for a specified
measurement model. A copy is measured with the binary POVM
$\{\Pi_z,I-\Pi_z\}$, where $\Pi_z=|z\rangle\langle z|$ is a normalized
multimode coherent-state projector. Classical adaptation of $z$ is allowed;
collective measurements, squeezing, and other measurement effects are not.
This is not the continuous-outcome heterodyne POVM. Unlike the earlier shell
injectivity theorem, the upper bounds below use points throughout a ball and
do not assume bounded state rank. A photon cutoff is used in the estimator,
with a proved bias bound, not imposed on the unknown state.

\subsection{State class, data metric, and main statements}
Fix a number of modes $s\ge1$, a moment exponent $\beta>0$, and $E>0$. Let
$\mathcal N=\sum_{j=1}^s a_j^*a_j$ be the total photon-number operator and set
\[
 \mathcal S_{s,\beta,E}
 =\{\rho\ge0:\operatorname{Tr}\rho=1,
                    \operatorname{Tr}(\rho\mathcal N^\beta)\le E\}.
\]
Thus $E$ bounds the moment itself; at $\beta=1$ it is the mean photon number.
Write $p_\rho(z)=\langle z|\rho|z\rangle=\pi^sQ_\rho(z)$ and
$D(\rho,\sigma)=\tfrac12\|\rho-\sigma\|_1$.
For $0<R\le\infty$ define
\[
 d_R(\rho,\sigma)=\sup_{\|z\|\le R}|p_\rho(z)-p_\sigma(z)|,
 \quad
 \Omega_E(R,\delta)=\sup_{\substack{\rho,\sigma\in\mathcal S_{s,\beta,E}\\
                                  d_R(\rho,\sigma)\le\delta}}
                                  D(\rho,\sigma).
\]
An infinite radius means all of $\mathbb C^s$. Use the conventions
$\log_+ x=\max(0,\log x)$ and $\log_+(x/\infty)=0$, and put
\[
 G_R(T)=1+\log_+(T/R^2),\qquad K_R(T)=T/G_R(T).
 \tag{A.1}\label{eq:aperture-scale}
\]
All asymptotic constants below may depend on the fixed $s,\beta,E,R_0$,
but are uniform for $R\ge R_0>0$, including radii that vary with $T$.
No uniform claim as the number of modes tends to infinity is made.

\begin{theorem}[Sharp deterministic modulus on a ball]
\label{thm:aperture-modulus}
For $T=\log(1/\delta)$ sufficiently large,
\[
 \boxed{\quad
 c\,K_R(T)^{-\beta/2}
 \le\Omega_E(R,\delta)\le C\,K_R(T)^{-\beta/2}.
 \quad}
 \tag{A.2}\label{eq:aperture-modulus}
\]
The lower bound already holds for pairs of pure states of moment exactly $E$.
The upper bound holds for states of arbitrary, possibly infinite rank.
\end{theorem}

For $B$ copies, let $\mathfrak R_B(R)$ be the infimum, over all adaptive
binary coherent-test protocols and density-operator estimators, of
$\sup_{\rho\in\mathcal S_{s,\beta,E}}\mathbb E_\rho D(\widehat\rho,\rho)$.
At each round the chosen amplitude has norm at most $R$ and a fresh copy is
measured. The protocol may depend on $B,s,\beta,E,R$.

\begin{theorem}[Sharp minimax rate, including classical adaptation]
\label{thm:aperture-minimax}
For all sufficiently large integers $B$,
\[
 \boxed{\quad
 c\,K_R(\log B)^{-\beta/2}
 \le\mathfrak R_B(R)\le C\,K_R(\log B)^{-\beta/2}.
 \quad}
 \tag{A.3}\label{eq:aperture-minimax}
\]
A nonadaptive tensor-Chebyshev design attains the upper order. The lower
bound holds even if the unknown state is promised pure, and even against
randomized adaptive choices of coherent amplitudes.
\end{theorem}

In particular, at fixed finite $R$ and $\beta=1$ the rate is
$\sqrt{\log\log B/\log B}$; with arbitrary amplitudes it is
$(\log B)^{-1/2}$. The transition is described by the whole factor $G_R$,
not just these endpoints. For instance $R^2=\log B/(\log\log B)^a$, with
fixed $a>0$, gives order $\sqrt{\log\log\log B/\log B}$ at $\beta=1$.
A bounded-rank promise with any rank bound at least one does not improve
these minimax orders, when the estimate may be an arbitrary density operator.
These are rates up to fixed multiplicative constants, not sharp leading
constants or claims about optimal runtime.

\subsection{A finite interpolation estimate}
Let $P_K$ project onto total occupations at most $K$, and let
$d_K=\binom{K+s}{s}$. For $0<S\le R^2$ put $a=\sqrt{S/(2s)}$.
For an integer $L\ge\max(2s,16S)$ use the $(L+1)^{2s}$ nodes whose real
coordinates are
\[
 x_j=a\cos\frac{(j+1/2)\pi}{L+1},\qquad 0\le j\le L.
 \tag{A.4}\label{eq:aperture-nodes}
\]
Each complex node lies in the ball of squared radius $S$.
Define
\begin{align}
 A_K(S)&=\max\{1,K!(2s/S)^K\},\nonumber\\
 C_{s,L}&=4^s(L+1)^{2s}3^L,\nonumber\\
 \tau_{s,L}(S)&=4\sqrt{\binom{L+2s}{2s-1}}
                    \frac{(2S)^{(L+1)/2}}{\sqrt{(L+1)!}}.
 \tag{A.5}\label{eq:aperture-constants}
\end{align}

\begin{lemma}[Low coefficients from a finite probability grid]
\label{lem:aperture-coefficients}
Suppose $L\ge2K$ and two states obey
$|p_\rho-p_\sigma|\le\eta$ at all nodes in \eqref{eq:aperture-nodes}.
Then, for all occupation multiindices $|\mu|,|\nu|\le K$,
\[
 |(\rho-\sigma)_{\mu\nu}|
 \le C_{s,L}A_K(S)\{e^S\eta+\tau_{s,L}(S)\}.
 \tag{A.6}\label{eq:aperture-coefficient}
\]
There is also an explicit linear reconstruction from noisy values of a single
state, with the same bound on each low matrix entry if the nodal probability
errors are at most $\eta$.
\end{lemma}
\begin{proof}
For $\Delta=\rho-\sigma$ define
\[
 F_\Delta(z)=e^{\|z\|^2}p_\Delta(z)
   =\sum_{\mu,\nu}\frac{\Delta_{\mu\nu}
                           \overline z^\mu z^\nu}{\sqrt{\mu!\nu!}}.
 \tag{A.7}\label{eq:aperture-entire}
\]
The entries of $\Delta$ have modulus at most two. For total degree $j$,
Cauchy--Schwarz and the multinomial identity bound the absolute sum on
$\|z\|^2\le S$ by
\[
 2\sqrt{\binom{j+2s-1}{2s-1}}\frac{(2S)^{j/2}}{\sqrt{j!}}.
\]
The ratio of consecutive summands without the leading two is
$\sqrt{2S(j+2s)}/(j+1)\le1/2$ for $j\ge L$ under the stated restrictions.
Consequently the total-degree-$L$ Taylor truncation $T_L$ has uniform error
at most $\tau_{s,L}(S)$. The same bound works for a single state.

In $2s$ real coordinates $T_L$ is a polynomial of total degree at most $L$.
Tensor discrete Chebyshev orthogonality at the chosen nodes recovers its
Chebyshev coefficients exactly. Keep only indices of total degree at most
$L$; all others are zero for $T_L$. The nodal error in $F_\Delta$ is bounded
by $e^S\eta+\tau_{s,L}(S)$. Each recovered coefficient error is at most
$4^s$ times this bound, because every tensor cosine has modulus at most one
and each nonconstant one-dimensional coefficient has prefactor two.

The sum of absolute monomial coefficients of the Chebyshev polynomial $T_j$
is at most $3^j$, by its recurrence. Replacing each real coordinate by
$(z_\ell+\overline z_\ell)/2$ or
$(z_\ell-\overline z_\ell)/(2i)$ does not increase the sum of absolute
coefficients of a homogeneous monomial. Scaling by $a$ contributes exactly
$a^{-d}$ to coefficients of total degree $d$. There are at most
$(L+1)^{2s}$ retained Chebyshev terms. It follows that the coefficient error
at $\overline z^\mu z^\nu$ is bounded by
$C_{s,L}a^{-|\mu|-|\nu|}(e^S\eta+\tau_{s,L})$.

Finally, $\mu!\le|\mu|!$. The sequence $\sqrt{j!}\,a^{-j}$ is log-convex,
so its maximum for $0\le j\le K$ is attained at an endpoint. Multiplying
two such maxima bounds
$\sqrt{\mu!\nu!}a^{-|\mu|-|\nu|}$ by $A_K(S)$.
The coefficients of $T_L$ below total degree $2K$ are the exact corresponding
coefficients of \eqref{eq:aperture-entire}. This proves the estimate and
describes the reconstruction. Hermitian symmetrization cannot enlarge the
largest entrywise error relative to the true Hermitian block.
\end{proof}

\begin{lemma}[State reconstruction and controlled cutoff bias]
\label{lem:aperture-state}
If the reconstructed Hermitian block $M$ has entrywise error at most $b$
relative to $P_K\rho P_K$, take its positive part and normalize its trace,
using the vacuum state if the positive part is zero. The resulting state
$\widehat\rho$, supported on $P_K$, satisfies
\[
 D(\widehat\rho,\rho)
 \le \tfrac32\sqrt{E/(K+1)^\beta}+2d_K^{3/2}b
 \tag{A.8}\label{eq:aperture-reconstruction}
\]
whenever $E/(K+1)^\beta\le1/2$. For two states in the moment class,
\[
 D(\rho,\sigma)
 \le2\sqrt{E/(K+1)^\beta}
    +\tfrac12d_K^{3/2}\max_{|\mu|,|\nu|\le K}|\Delta_{\mu\nu}|.
 \tag{A.9}\label{eq:aperture-two-state}
\]
\end{lemma}
\begin{proof}
The omitted probability $q=\operatorname{Tr}((I-P_K)\rho)$ is at most
$E/(K+1)^\beta$. Decomposing $\rho$ into rank-one states and applying
Cauchy--Schwarz gives
$\|\rho-P_K\rho P_K\|_1\le2\sqrt q$.
Entrywise error $b$ implies trace-norm block error
$h\le d_K^{3/2}b$ by the Hilbert--Schmidt inequality.
If $A=M_+$, positivity of the true block implies
$\operatorname{Tr}M_-\le h$, hence
$\|A-P_K\rho P_K\|_1\le2h$.
If $A\ne0$, normalizing it changes trace norm by
$|1-\operatorname{Tr}A|\le q+2h$.
Thus $D\le\sqrt q+q/2+2h$, which proves \eqref{eq:aperture-reconstruction}.
If $A=0$, then $h\ge1-q\ge1/2$, so the displayed bound is at least one
and remains valid for the vacuum fallback. Applying the cutoff bound to
both states and the block trace-norm bound proves \eqref{eq:aperture-two-state}.
\end{proof}

\subsection{Balancing the aperture and the cutoff}
Set
\[
 \Phi_R(K)=K\{1+\log_+(K/R^2)\},\quad
 S=\min(R^2,K),\quad L=128sK.
 \tag{A.10}\label{eq:aperture-phi}
\]
There is a constant $c_s$ such that the noise multiplier in
\eqref{eq:aperture-coefficient} obeys
\[
 C_{s,L}A_K(S)e^S\le C_sK^{2s}\exp(c_s\Phi_R(K)).
 \tag{A.11}\label{eq:aperture-noise}
\]
This follows from $K!\le K^K$, $S\le K$, and
$\log(K/S)=\log_+(K/R^2)$.
For the bias term one has the stronger uniform estimate
\[
 C_{s,L}A_K(S)\tau_{s,L}(S)
 \le C_sK^{3s-1/2}2^{-(64s-1)K}.
 \tag{A.12}\label{eq:aperture-bias}
\]
Indeed, use $n!\ge(n/e)^n$ in \eqref{eq:aperture-constants}. Apart from
the displayed polynomial factors, the left side is at most
\[
 (2s)^K\left(\frac{18e}{128s}\right)^{64sK}
        (S/K)^{(64s-1)K}
 \le(2s)^{-(64s-1)K}\le2^{-(64s-1)K},
\]
since $e<3$ and $54/128<1/2$. The deliberately large $128s$ is a safe
analytic choice, not an optimized numerical parameter.

The elementary inverse-scale estimate is
\[
 (1-e^{-1})T\le\Phi_R(K_R(T))\le T.
 \tag{A.13}\label{eq:aperture-inversion}
\]
For $T/R^2\le1$ it is immediate. Otherwise set
$g=1+\log(T/R^2)\ge1$. Then $K_R(T)/R^2\ge1$ and the middle quantity equals
$T(1-\log g/g)$; the maximum of $\log g/g$ on $g\ge1$ is $e^{-1}$.
Also $\Phi_R(cK)\le c\Phi_R(K)$ for $0<c\le1$, and the reversed inequality
holds for $c\ge1$. Integer rounding changes the resulting large-$T$ scales
only by fixed factors, uniformly for $R\ge R_0>0$.

To obtain the deterministic upper bound, take
$K=\lfloor cK_R(T)\rfloor$ with a sufficiently small fixed $c>0$.
Equations \eqref{eq:aperture-noise}--\eqref{eq:aperture-inversion} make the
noise contribution, with $\eta=\delta=e^{-T}$, at most a polynomial in $T$
times $e^{-T/2}$. The bias is exponentially small in $K$ by
\eqref{eq:aperture-bias}. Equation \eqref{eq:aperture-two-state} gives
$C K^{-\beta/2}$, as claimed. Only a finite subset of the ball was used.

\subsection{The two-level lower bound, uniformly over amplitudes}
For large integers $n$ let $u=E/n^\beta\le1/2$ and use states
\[
 |\psi_n^\pm\rangle=\sqrt{1-u}|0,\ldots,0\rangle
       \pm\sqrt u|n,0,\ldots,0\rangle,\qquad
 \rho_n^\pm=|\psi_n^\pm\rangle\langle\psi_n^\pm|.
 \tag{A.14}\label{eq:aperture-pair}
\]
Their moment is exactly $E$ and
\[
 D(\rho_n^+,\rho_n^-)=2\sqrt{u(1-u)}\ge\sqrt{2E}\,n^{-\beta/2}.
 \tag{A.15}\label{eq:aperture-pair-distance}
\]
Their coherent amplitudes are the sum and difference of two terms of moduli
\[
 a_z=\sqrt{1-u}\,e^{-\|z\|^2/2},\qquad
 b_z=\sqrt u\,e^{-\|z\|^2/2}\frac{|z_1|^n}{\sqrt{n!}}.
\]
Their relative phase is $n\arg z_1$. Maximizing $e^{-x}x^{n/2}$ gives
\[
 |p_+(z)-p_-(z)|\le4a_zb_z,
 \quad a_zb_z\le\sqrt u\,2^{-n/2},
 \quad b_z^2\le u/2.
 \tag{A.16}\label{eq:aperture-amplitudes}
\]
Here $e^{-n}n^n/n!\le1/e<1/2$ for $n\ge1$; the sequence is decreasing
because $(1+1/n)^n<e$. Other modes only decrease both amplitudes.
It follows that
\[
 d_R(\rho_n^+,\rho_n^-)
 \le4\sqrt u\min\left\{2^{-n/2},\frac{R^n}{\sqrt{n!}}\right\}.
 \tag{A.17}\label{eq:aperture-forward-small}
\]
Using $n!\ge(n/e)^n$, there is an absolute $c_0>0$ such that the right side
is at most $4\exp(-c_0\Phi_R(n))$. To check uniformity, write
$t=\log_+(n/R^2)$. The two exponents are
$(\log2)n/2$ and $(t-1)n/2$; their maximum bounds a positive constant times
$(1+t)n$. Choose $n=\lceil C K_R(T)\rceil$ with sufficiently large fixed $C$.
Then \eqref{eq:aperture-inversion} makes \eqref{eq:aperture-forward-small}
at most $e^{-T}$ for large $T$, and \eqref{eq:aperture-pair-distance} proves
the lower bound in Theorem~\ref{thm:aperture-modulus}.

\subsection{Adaptive experiments cannot beat the lower rate}
For Bernoulli parameters $p_+,p_-$ define the unnormalized squared Hellinger
distance
\[
 h^2(p_+,p_-)=(\sqrt{p_+}-\sqrt{p_-})^2
                  +(\sqrt{1-p_+}-\sqrt{1-p_-})^2.
\]
For all large $n$, uniformly in $z\in\mathbb C^s$, one has
$p_\pm(z)\le1-u/4$. Indeed \eqref{eq:aperture-amplitudes} bounds these
probabilities by $1-u+u/2+2\sqrt u\,2^{-n/2}$, and the last term is at most
$u/4$ for large $n$, since $u=E/n^\beta$.
The difference of the square roots of the first probabilities is at most
$2\min(a_z,b_z)$. For the complementary probabilities, rationalizing their
square-root difference and using $1-p_\pm\ge u/4$ gives an upper bound
$64a_z^2b_z^2/u$ for its square. Therefore
\begin{align}
 h^2(p_+(z),p_-(z))
 &\le4\min(a_z^2,b_z^2)+64a_z^2b_z^2/u,\nonumber\\
 \sup_{\|z\|\le R}h^2(p_+(z),p_-(z))
 &\le68\min\left\{2^{-n/2},\frac{R^{2n}}{n!}\right\}
 \le68e^{-c_1\Phi_R(n)}
 \tag{A.18}\label{eq:aperture-hellinger}
\end{align}
for an absolute $c_1>0$. For the first of the two uniform bounds use
$\min(a_z^2,b_z^2)\le a_zb_z$; for the second use
$\min(a_z^2,b_z^2)\le b_z^2$ and drop Gaussian factors.
The last inequality follows by comparing
$(\log2)/2$ and $t-1$ for $t=\log_+(n/R^2)$.

For completeness, the same estimate controls an adaptive transcript.
Conditional on any common history the protocol uses the same distribution
of actions under the two hypotheses. Its conditional Hellinger affinity is
at least $1-\eta/2$ when the bound in \eqref{eq:aperture-hellinger} is $\eta$.
Backward integration through the $B$ rounds shows total affinity at least
$(1-\eta/2)^B\ge1-B\eta/2$. Thus transcript Hellinger distance squared is
at most $B\eta$, and total variation is at most $\sqrt{B\eta}$.
This argument includes randomized actions by conditioning also on the random
seed, and does not suppose nonadaptive or independent action choices.

Take $n=\lceil C K_R(\log B)\rceil$ with $C$ large enough that
$\eta\le68B^{-3}$. For large $B$ the transcript total variation is at most
$1/2$. Given any state estimate, classify it by the closer of
$\rho_n^+$ and $\rho_n^-$. A classification error forces trace-distance
loss at least half the distance in \eqref{eq:aperture-pair-distance}.
The two-point testing bound therefore gives maximum expected loss at least
one eighth of that distance. This proves the adaptive minimax lower bound.

\subsection{A nonadaptive estimator attains the rate}
Set $T=\log B$ and $K=\lfloor cK_R(T)\rfloor$, with $c>0$ chosen sufficiently
small depending on $s$. Use \eqref{eq:aperture-phi} and repeat every grid
setting $b=\lfloor B/(L+1)^{2s}\rfloor$ times. Unused copies may be discarded.
Hoeffding's bound and a union bound show that, with probability at least
$1-B^{-2}$, every empirical probability has error at most
\[
 \eta_B=\sqrt{\frac{\log(2(L+1)^{2s}B^2)}{2b}}
            =\exp(-T/2+O_s(\log T)).
 \tag{A.19}\label{eq:aperture-hoeffding}
\]
This concentration follows directly from the Bernoulli moment-generating
function bound $\mathbb E e^{t(X-\mathbb EX)}\le e^{t^2/8}$ and Markov's
inequality; no assumption on the unknown rank is used.
By decreasing $c$, \eqref{eq:aperture-noise} amplifies this by at most
$e^{T/4}$ up to polynomial factors. Use the linear reconstruction in
Lemma~\ref{lem:aperture-coefficients} and the positive-part estimator in
Lemma~\ref{lem:aperture-state}. On the concentration event its error is
at most $C K^{-\beta/2}$, since the analytic bias remains exponentially
small by \eqref{eq:aperture-bias}. On the complement its error is at most one.
This proves the expected upper bound, completing
Theorem~\ref{thm:aperture-minimax}.

The estimator has a finite classical output matrix of dimension $d_K$.
Its grid has $(128sK+1)^{2s}$ settings, which is polynomial in $\log B$ for
fixed $s$ but not dimension-independent. Large interpolation condition
numbers require sufficient arithmetic precision. Rounding settings to within
$\epsilon$ perturbs a probability by at most $\epsilon$, because
$\|\Pi_z-\Pi_w\|_{\rm op}=\sqrt{1-e^{-\|z-w\|^2}}\le\|z-w\|$.
Thus logarithmically many precision digits in $B$ suffice for the asymptotic
scheme. This does not make the exponential copy requirement practical.

\subsection{Copy-complexity and measurement-model consequences}
Let $B_*(\varepsilon,R)$ be the least budget with minimax expected error at
most $\varepsilon$. For fixed $s,\beta,E,R_0$ and $\varepsilon\downarrow0$,
\[
 \boxed{
 \log B_*(\varepsilon,R)
 \asymp \varepsilon^{-2/\beta}
                \{1+\log_+(\varepsilon^{-2/\beta}/R^2)\},
 \qquad R\ge R_0.
 }
 \tag{A.20}\label{eq:aperture-complexity}
\]
This is equivalent to \eqref{eq:aperture-minimax} via
\eqref{eq:aperture-inversion}; constants are not uniform in $E$ or $\beta$.
For a mean-energy bound and fixed finite $R$, the copy complexity is
$\exp(\Theta(\varepsilon^{-2}\log(1/\varepsilon)))$; for arbitrary coherent
amplitudes it is $\exp(\Theta(\varepsilon^{-2}))$.
Classical adaptation does not change either order. Rank one already forces
the lower bounds.

This is a restriction imposed by the allowed measurements, not by quantum
mechanics alone. With arbitrary binary projectors a polynomial-copy scheme
at fixed $s,\beta,E$ is elementary. Choose
$K=O(\varepsilon^{-2/\beta})$ to control the cutoff bias, measure all
occupation projectors and the projectors onto
$(|j\rangle+|k\rangle)/\sqrt2$ and
$(|j\rangle+i|k\rangle)/\sqrt2$ within $P_K$.
Estimating these $O(d_K^2)$ probabilities to
$O(\varepsilon/d_K^{3/2})$ recovers the low block in trace norm and needs
$O(d_K^5\varepsilon^{-2}\log d_K)$ copies for fixed success probability.
Normalization as above completes the estimate. This deliberately nonoptimal
construction is sufficient to demonstrate a separation from coherent tests.
Sharper general continuous-variable learning bounds are prior work
\cite{Mele2025}. The coherent-test lower bound must not be applied to these
other projectors, to homodyne measurements, to collective protocols, or to
all Gaussian measurement schemes.

Interpolation-based continuous-variable tomography is established prior work
\cite{LGC2018}; the proposed contribution here is the matching trace-distance
minimax law with an aperture-dependent radius and a sole photon-moment
constraint, including adaptive lower bounds and the fixed-mode extension.
Exact injectivity from finitely many complete shells is a different theorem.
In particular, \eqref{eq:aperture-modulus} does not supply a matching upper
rate for the fixed-shell modulus in \eqref{eq:husimi-modulus}.
These results neither prove nor assume microscopic BGS universality.
External priority and independent referee acceptance remain unestablished.
